\documentclass[a4paper,12pt]{article}

\usepackage[T2A]{fontenc}
\usepackage[utf8]{inputenc}
\usepackage[russian]{babel}
\usepackage{amsmath,amssymb,mathtools}
\usepackage{PTSans}
\renewcommand{\familydefault}{\sfdefault}
\usepackage{icomma}
\usepackage[a4paper,margin=20mm]{geometry}
\usepackage{microtype}
\usepackage{tikz}
\usetikzlibrary{calc,arrows.meta,positioning,shapes.geometric}
\usepackage{tabularx,booktabs,longtable}
\usepackage{enumitem}
\usepackage{hyperref}
\usepackage{parskip}
\usepackage{fancyhdr}
\usepackage{setspace}
\usepackage{xcolor}

\definecolor{accent}{HTML}{287A78}
\definecolor{darkaccent}{HTML}{174F4D}
\definecolor{textgray}{HTML}{303030}
\definecolor{softgray}{HTML}{F2F4F4}
\definecolor{muted}{HTML}{687676}

\hypersetup{colorlinks=true,linkcolor=darkaccent,urlcolor=darkaccent}
\setlength{\headheight}{25pt}
\onehalfspacing
\setlength{\parindent}{0pt}
\setlist[itemize]{leftmargin=1.4em,itemsep=0.28em,topsep=0.35em}
\setlist[enumerate]{leftmargin=1.7em,itemsep=0.45em,topsep=0.4em}
\pagestyle{fancy}
\fancyhf{}
\fancyhead[L]{\small\color{muted} Рациональные неравенства}
\fancyhead[R]{\small\color{muted} Анна Трапезникова}
\fancyfoot[C]{\small\color{muted}\thepage}
\renewcommand{\headrulewidth}{0.3pt}
\renewcommand{\headrule}{\hbox to\headwidth{\color{accent}\leaders\hrule height \headrulewidth\hfill}}

\newcommand{\checkitem}{\(\square\)\hspace{0.45em}}
\newcommand{\important}[1]{\textcolor{darkaccent}{\textbf{#1}}}
\newcommand{\R}{\mathbb{R}}

\begin{document}

% ---------- TITLE ----------
\thispagestyle{empty}
\begin{center}
\vspace*{24mm}
{\color{accent}\rule{0.72\linewidth}{1.2pt}}\\[10mm]
{\Huge\bfseries\color{darkaccent} Чек-лист}\\[7mm]
{\LARGE\bfseries Рациональные неравенства}\\[3mm]
{\Large Метод интервалов, кратность корней и нормализация}\\[14mm]
{\large 27.09.26}\\[22mm]

\begin{minipage}{0.76\linewidth}
\centering
\large
Главная цель: уверенно приводить рациональное неравенство к удобному виду, находить все критические точки, учитывать кратность корней и безошибочно выбирать нужные промежутки.
\end{minipage}

\vfill
{\large\bfseries Лёвин Артём Александрович}\\[2mm]
{\normalsize эксклюзивно для Анны Трапезниковой}
\vspace{13mm}

\begin{tikzpicture}[remember picture,overlay]
  \draw[accent,line width=1.4pt]
    ($(current page.south west)+(18mm,18mm)$)
    .. controls ($(current page.south west)+(60mm,29mm)$) and ($(current page.south east)+(-64mm,7mm)$)
    .. ($(current page.south east)+(-18mm,20mm)$);
  \draw[accent!55,line width=0.75pt]
    ($(current page.south west)+(18mm,13mm)$)
    .. controls ($(current page.south west)+(75mm,3mm)$) and ($(current page.south east)+(-78mm,27mm)$)
    .. ($(current page.south east)+(-18mm,14mm)$);
\end{tikzpicture}
\end{center}
\clearpage

% ---------- CORE CHECKLIST ----------
\section*{1. Что должно быть доведено до автоматизма}

\textbf{Нормальная форма рационального неравенства:}
\[
\frac{P(x)}{Q(x)}\;\gtreqless\;0, \qquad Q(x)\ne0.
\]
Именно к этому виду удобно применять метод интервалов. Если справа стоит дробь или другое выражение, сначала перенесите всё в одну сторону и приведите к общему знаменателю.

\begin{itemize}
  \item \checkitem \textbf{ОДЗ записано до сокращений.} Корни знаменателя никогда не входят в ответ.
  \item \checkitem \textbf{Числитель и знаменатель разложены на множители.} Не раскрывайте знаменатель без необходимости: факторизованный вид полезнее.
  \item \checkitem \textbf{Найдены все критические точки.} Это нули числителя и нули знаменателя.
  \item \checkitem \textbf{Для каждой точки определена кратность.} Степень множителя \((x-a)^k\) и есть кратность корня \(a\).
  \item \checkitem \textbf{Точки нанесены на числовую прямую в правильном порядке.}
  \item \checkitem \textbf{Один знак найден подстановкой.} Достаточно определить знак выражения, вычислять точное значение обычно не нужно.
  \item \checkitem \textbf{Учтено правило кратности.} Через корень нечётной кратности знак меняется, через корень чётной кратности --- сохраняется.
  \item \checkitem \textbf{Ответ собран по знаку неравенства.} Нули числителя включаются только при \(\le\) или \(\ge\); нули знаменателя исключаются всегда.
\end{itemize}

\subsection*{Кратность: главное правило}
Пусть в разложении есть множитель \((x-a)^k\).
\begin{itemize}
  \item Если \(k\) \textbf{нечётно}, выражение меняет знак при переходе через \(a\).
  \item Если \(k\) \textbf{чётно}, выражение сохраняет знак при переходе через \(a\).
\end{itemize}

\important{Важно:} чётность относится к \emph{кратности корня}, то есть к количеству повторений соответствующего множителя, а не к самому числу \(a\).

\subsection*{Как отмечать точки}
\begin{center}
\renewcommand{\arraystretch}{1.25}
\begin{tabularx}{0.94\linewidth}{@{}X X X@{}}
\toprule
\textbf{Точка} & \textbf{Строгое неравенство} & \textbf{Нестрогое неравенство} \\
\midrule
Ноль числителя & исключить & включить \\
Ноль знаменателя & исключить & исключить \\
\bottomrule
\end{tabularx}
\end{center}

% ---------- WORKED EXAMPLE 1 ----------
\section*{2. Разобранный пример: кратности корней}
Решим
\[
\frac{(x-1)^2(x+3)^6(x-5)^9}{x^2}\ge0.
\]

\textbf{Шаг 1. Критические точки.}
\[
\begin{array}{c|c|c}
\text{точка} & \text{откуда} & \text{кратность}\\
\hline
-3 & (x+3)^6 & 6\ \text{(чётная)}\\
0 & x^2\ \text{в знаменателе} & 2\ \text{(чётная)}\\
1 & (x-1)^2 & 2\ \text{(чётная)}\\
5 & (x-5)^9 & 9\ \text{(нечётная)}
\end{array}
\]
Точка \(0\) запрещена; точки \(-3,1,5\) могут входить в ответ, потому что неравенство нестрогое.

\textbf{Шаг 2. Один проверочный знак.} Возьмём \(x=2\):
\[
(2-1)^2>0,\quad (2+3)^6>0,\quad (2-5)^9<0,\quad 2^2>0,
\]
значит, на интервале \((1,5)\) стоит знак «\(-\)».

\textbf{Шаг 3. Распространяем знаки.} Через \(-3\), \(0\), \(1\) знак не меняется; через \(5\) меняется.

\begin{center}
\begin{tikzpicture}[x=1.45cm,y=1cm,>=Latex]
  \draw[->,line width=0.8pt] (-4.3,0) -- (6.4,0) node[right] {$x$};
  \foreach \x/\lab/\kind in {-3/-3/fill,0/0/open,1/1/fill,5/5/fill}{
    \draw (\x,0.09)--(\x,-0.09);
    \node[below=3pt] at (\x,0) {$\lab$};
  }
  \fill[accent] (-3,0) circle (2.2pt);
  \draw[accent,line width=1pt,fill=white] (0,0) circle (2.4pt);
  \fill[accent] (1,0) circle (2.2pt);
  \fill[accent] (5,0) circle (2.2pt);
  \node[above=5pt] at (-3.65,0) {$-$};
  \node[above=5pt] at (-1.5,0) {$-$};
  \node[above=5pt] at (0.5,0) {$-$};
  \node[above=5pt] at (3,0) {$-$};
  \node[above=5pt] at (5.65,0) {$+$};
\end{tikzpicture}
\end{center}

\[
\boxed{x\in\{-3;1\}\cup[5;+\infty)}.
\]

\subsection*{Что здесь проверяется}
\begin{itemize}
  \item чётная кратность сохраняет знак;
  \item нечётная кратность меняет знак;
  \item запрещённая точка остаётся выколотой независимо от кратности;
  \item при нестрогом знаке нули числителя включаются.
\end{itemize}

\subsection*{Типичная ошибка}
Не ставьте знаки «по памяти» сразу на всех интервалах. Сначала найдите \textbf{один} знак подстановкой, затем переходите через точки с учётом кратностей.

% ---------- WORKED EXAMPLE 2 ----------
\section*{3. Разобранный пример: сначала нормализовать}
Решим
\[
\frac{3}{x-5}<\frac{3x-1}{x^2-1}.
\]

\textbf{1. Переносим всё влево и приводим к общему знаменателю:}
\begin{align*}
\frac{3}{x-5}-\frac{3x-1}{x^2-1}&<0,\\
\frac{3(x^2-1)-(3x-1)(x-5)}{(x-5)(x^2-1)}&<0,\\
\frac{16x-8}{(x-5)(x-1)(x+1)}&<0,\\
\frac{8(2x-1)}{(x-5)(x-1)(x+1)}&<0.
\end{align*}

\important{Не раскрывайте знаменатель:} его факторизованный вид сразу показывает запрещённые точки.

\textbf{2. Критические точки:}
\[
-1,\quad \frac12,\quad 1,\quad 5.
\]
Все имеют нечётную кратность, поэтому при проходе через каждую точку знак меняется. Точки \(-1,1,5\) запрещены; \(\frac12\) тоже не входит, потому что неравенство строгое.

\textbf{3. Ответ:}
\[
\boxed{x\in\left(-1;\frac12\right)\cup(1;5)}.
\]

\subsection*{Когда можно сокращать общий множитель}
Если числитель и знаменатель имеют общий множитель, сначала зафиксируйте ограничение из знаменателя.
Например,
\[
\frac{(x-3)(x+3)}{(x-3)(x+2)}\ge0,\qquad x\ne3,-2.
\]
После сокращения можно исследовать \(\dfrac{x+3}{x+2}\), но точка \(x=3\) остаётся запрещённой и должна быть выколота в финальном ответе.

\clearpage

% ---------- FLOWCHART ----------
\thispagestyle{fancy}
\begin{center}
{\LARGE\bfseries\color{darkaccent} Алгоритм: рациональное неравенство методом интервалов}
\end{center}
\vspace{5mm}

\begin{center}
\begin{tikzpicture}[
  node distance=6.5mm,
  startstop/.style={ellipse,draw=accent,line width=1pt,minimum width=39mm,minimum height=10mm,align=center,fill=softgray},
  process/.style={rectangle,rounded corners=2mm,draw=accent,line width=1pt,text width=125mm,minimum height=10mm,align=center,inner sep=5pt},
  decision/.style={diamond,aspect=2.4,draw=accent,line width=1pt,text width=36mm,align=center,inner sep=2pt},
  line/.style={-{Latex[length=3mm]},line width=0.9pt,draw=darkaccent}
]
  \node[startstop] (start) {Старт};
  \node[process,below=of start] (norm) {Перенесите всё в одну сторону и получите \(\dfrac{P(x)}{Q(x)}\gtreqless0\)};
  \node[process,below=of norm] (domain) {Запишите ограничения: \(Q(x)\ne0\)};
  \node[process,below=of domain] (factor) {Разложите числитель и знаменатель на множители; знаменатель без нужды не раскрывайте};
  \node[process,below=of factor] (points) {Найдите нули числителя и знаменателя, определите кратности, расположите точки по возрастанию};
  \node[process,below=of points] (sign) {Подставьте одно удобное число и определите знак на одном интервале};
  \node[decision,below=8mm of sign] (mult) {Переходим через очередную критическую точку};
  \node[process,below=8mm of mult,xshift=-31mm,text width=48mm] (odd) {Нечётная кратность: знак меняется};
  \node[process,below=8mm of mult,xshift=31mm,text width=48mm] (even) {Чётная кратность: знак сохраняется};
  \node[process,below=25mm of mult,text width=116mm] (select) {Выберите интервалы нужного знака. Нули числителя включайте только при \(\le\) или \(\ge\); нули знаменателя исключайте всегда};
  \node[startstop,below=of select] (finish) {Запишите ответ и проверьте ОДЗ};

  \draw[line] (start)--(norm);
  \draw[line] (norm)--(domain);
  \draw[line] (domain)--(factor);
  \draw[line] (factor)--(points);
  \draw[line] (points)--(sign);
  \draw[line] (sign)--(mult);
  \draw[line] (mult)--node[above left,sloped]{\small нечётная}(odd);
  \draw[line] (mult)--node[above right,sloped]{\small чётная}(even);
  \draw[line] (odd.south) |- (select.west);
  \draw[line] (even.south) |- (select.east);
  \draw[line] (select)--(finish);
\end{tikzpicture}
\end{center}
\clearpage

% ---------- SELF-CHECK ----------
\section*{4. Быстрая самопроверка перед ответом}
Перед тем как закончить решение, задайте себе семь вопросов.
\begin{enumerate}
  \item Приведено ли неравенство к виду «одна дробь и ноль»?
  \item Выписаны ли все запрещённые значения знаменателя?
  \item Все ли выражения разложены на множители настолько, насколько это полезно?
  \item Не потеряна ли запрещённая точка после сокращения общего множителя?
  \item Для всех критических точек верно определена кратность?
  \item Проверен ли хотя бы один знак реальной подстановкой?
  \item Включены ли только допустимые граничные точки?
\end{enumerate}

\subsection*{Пять типичных ошибок}
\begin{itemize}
  \item Умножать обе части на знаменатель неизвестного знака и забывать, что знак неравенства может измениться.
  \item Сократить общий множитель и забыть выколоть его ноль из знаменателя.
  \item Менять знак при переходе через корень чётной кратности.
  \item Включать ноль знаменателя при нестрогом неравенстве.
  \item Раскрывать факторизованный знаменатель и тем самым усложнять поиск критических точек.
\end{itemize}

\important{Ориентир:} в методе интервалов Вы исследуете прежде всего \textbf{знак} выражения. Точное числовое значение при подстановке обычно не требуется.

Если множитель заведомо положителен, его можно учесть сразу. Например, \(x^4+x^2+1>0\) при любом \(x\in\R\): он не создаёт критических точек и не меняет знак дроби.

\clearpage

% ---------- TRAINING ----------
\section*{5. Тренировочный блок — Путь А}
\textbf{10 задач.} Решите методом интервалов. В каждой задаче явно отметьте ОДЗ и критические точки. Решений здесь нет; ответы приведены на следующей странице.

\begin{enumerate}
  \item \(\displaystyle \frac{(x-2)(x+3)}{x-1}\ge0.\)

  \item \(\displaystyle \frac{(x+1)^2(x-4)}{x-2}>0.\)

  \item \(\displaystyle \frac{(x-3)^2(x+2)^3}{(x+1)^2}\le0.\)

  \item \(\displaystyle \frac{x^2-9}{x^2-x-6}\ge0.\)

  \item \(\displaystyle \frac{3}{x-5}<\frac{3x-1}{x^2-1}.\)

  \item \(\displaystyle \frac{x^2-4x}{x^2-5x+6}\ge0.\)

  \item \(\displaystyle \frac{x^2(x-4)^2}{(x-1)(x-5)}\ge0.\)

  \item \(\displaystyle \frac{(x+1)^4(x-2)^3}{(2x-5)^2(x+4)}<0.\)

  \item \(\displaystyle \frac{x^4+x^2+1}{(x+1)(x-5)}\le0.\)

  \item \(\displaystyle \frac{x^4-5x^2+4}{x^2-4}\ge0.\)
\end{enumerate}

\vfill
\begin{center}
\small\color{muted}
Проверяйте не только интервалы, но и судьбу каждой граничной точки.
\end{center}
\clearpage

% ---------- ANSWERS ----------
\section*{6. Ответы к тренировочному блоку}

\renewcommand{\arraystretch}{1.5}
\begin{longtable}{@{}p{0.08\linewidth}p{0.84\linewidth}@{}}
\toprule
\textbf{№} & \textbf{Ответ}\\
\midrule
1 & \(\displaystyle [-3;1)\cup[2;+\infty)\)\\
2 & \(\displaystyle (-\infty;-1)\cup(-1;2)\cup(4;+\infty)\)\\
3 & \(\displaystyle (-\infty;-2]\cup\{3\}\)\\
4 & \(\displaystyle (-\infty;-3]\cup(-2;3)\cup(3;+\infty)\)\\
5 & \(\displaystyle \left(-1;\frac12\right)\cup(1;5)\)\\
6 & \(\displaystyle (-\infty;0]\cup(2;3)\cup[4;+\infty)\)\\
7 & \(\displaystyle (-\infty;1)\cup\{4\}\cup(5;+\infty)\)\\
8 & \(\displaystyle (-4;-1)\cup(-1;2)\)\\
9 & \(\displaystyle (-1;5)\)\\
10 & \(\displaystyle (-\infty;-2)\cup(-2;-1]\cup[1;2)\cup(2;+\infty)\)\\
\bottomrule
\end{longtable}

\subsection*{Контрольные акценты}
\begin{itemize}
  \item В №4 точка \(x=3\) остаётся запрещённой после сокращения общего множителя.
  \item В №7 точка \(x=4\) даёт изолированное допустимое значение, потому что числитель обращается в ноль, а соседние интервалы отрицательны.
  \item В №9 числитель \(x^4+x^2+1\) всегда положителен, поэтому знак определяется только знаменателем.
  \item В №10 после сокращения множителя \(x^2-4\) необходимо сохранить запрет \(x\ne\pm2\).
\end{itemize}

\vfill
\begin{center}
\color{darkaccent}\rule{0.58\linewidth}{0.7pt}\\[3mm]
\small Лёвин Артём Александрович эксклюзивно для Анны Трапезниковой
\end{center}

\end{document}
